\documentclass[12pt,a4paper,openany,oneside]{book}

% --- Paquetes ---
\usepackage[utf8]{inputenc}
\usepackage[spanish, es-noshorthands]{babel}
\usepackage{amsmath, amsfonts, amssymb}
\usepackage{graphicx}
\usepackage{geometry}
\usepackage{hyperref}
\usepackage{booktabs}
\usepackage{fancyhdr}
\usepackage{listings}
\usepackage{xcolor}
\usepackage{tikz}
\usepackage{pgfplots}
\usepackage{colortbl}
\usepackage{multirow}
\usepackage{hhline}
\usetikzlibrary{arrows.meta, positioning, shapes.geometric, calc}
\pgfplotsset{compat=1.18}

\geometry{margin=2.5cm}

% --- Colores y estilos para código Python ---
\definecolor{codegreen}{rgb}{0,0.6,0}
\definecolor{codegray}{rgb}{0.5,0.5,0.5}
\definecolor{codepurple}{rgb}{0.58,0,0.82}
\definecolor{backcolour}{rgb}{0.95,0.95,0.92}

\lstset{
    language=Python,
    backgroundcolor=\color{backcolour},   
    commentstyle=\color{codegreen},
    keywordstyle=\color{blue},
    numberstyle=\tiny\color{codegray},
    stringstyle=\color{codepurple},
    basicstyle=\ttfamily\small,
    breaklines=true,
    frame=single,
    showstringspaces=false
}

% --- Estilo de cabecera y pie ---
\pagestyle{fancy}
\fancyhf{}
\renewcommand{\headrulewidth}{0.4pt}
\renewcommand{\footrulewidth}{0.4pt}
\fancyhead[L]{\small Carlos César Sánchez Coronel}
\fancyhead[R]{\small Sesión 2: EDA y Feature Engineering Avanzado}
\fancyfoot[L]{\small Carlos César Sánchez Coronel}
\fancyfoot[C]{\thepage}

% --- Portada ---
\title{
    \textbf{Sesión 2} \\ 
    \Huge \textbf{ANÁLISIS EXPLORATORIO Y FEATURE ENGINEERING AVANZADO} \\
    \Large De los datos crudos a características de alto impacto
}
\author{
    \small Por \\ \vspace{0.2cm}
    \Large \textsc{Carlos César Sánchez Coronel}
}
\date{2026}

\begin{document}

\maketitle
\tableofcontents
\addcontentsline{toc}{chapter}{Índice general}

% ------------------------------------------------------------------
% CAPÍTULO 2: SESIÓN 2 - EDA Y FEATURE ENGINEERING AVANZADO
% ------------------------------------------------------------------
\chapter{Análisis Exploratorio y Feature Engineering Avanzado}

En esta sesión se aborda el proceso detectivesco que precede a cualquier modelado: el análisis exploratorio de datos (EDA) y la ingeniería de características (feature engineering). Se presentan técnicas avanzadas, sus fundamentos matemáticos y criterios de aplicación para transformar datos crudos en representaciones que maximicen el rendimiento de los modelos.

\section{Logro de la sesión}
Realizar un análisis exploratorio de datos profundo y aplicar técnicas avanzadas de feature engineering para transformar datos crudos en representaciones óptimas que potencien el rendimiento de modelos de machine learning, comprendiendo el fundamento matemático de cada técnica y su impacto en la calidad del modelado.

\section{Introducción: EDA como proceso detectivesco}
El análisis exploratorio de datos no es un mero requisito burocrático; es el proceso mediante el cual el científico de datos se convierte en detective, buscando pistas, anomalías y patrones que guiarán todas las decisiones posteriores. John Tukey, pionero del EDA, lo definió como una actitud filosófica y un conjunto de técnicas que permiten ``escuchar a los datos''. El feature engineering, por su parte, es el arte de traducir esos hallazgos en lenguaje matemático que los algoritmos puedan aprovechar. Juntos, EDA y feature engineering determinan el límite superior del rendimiento que cualquier modelo podrá alcanzar.

\section{Análisis Exploratorio de Datos Avanzado}

\subsection{Estadísticas descriptivas robustas vs. clásicas}
Las estadísticas clásicas (media, desviación estándar) son sensibles a outliers. Una sola observación extrema puede distorsionar completamente la media. Para conjuntos de datos reales (siempre contaminados), se requieren estimadores robustos.

\begin{table}[h]
\centering
\begin{tabular}{|l|l|l|l|}
\hline
\textbf{Medida} & \textbf{Clásica} & \textbf{Robusta} & \textbf{Fórmula robusta} \\
\hline
Tendencia central & $\bar{x} = \frac{1}{n}\sum x_i$ & Mediana & $Mediana = x_{(n/2)}$ \\
Dispersión & $s = \sqrt{\frac{1}{n-1}\sum(x_i-\bar{x})^2}$ & MAD & $MAD = mediana(|x_i - mediana(x)|)$ \\
Posición & Cuartiles & Cuartiles & $Q1, Q3$ (ya robustos) \\
\hline
\end{tabular}
\caption{Comparación entre estadísticos clásicos y robustos.}
\end{table}

\textbf{Propósito:} Detectar asimetrías y outliers comparando media y mediana. Si difieren significativamente, hay asimetría o valores extremos. Usar MAD en lugar de desviación estándar para detectar outliers mediante el z-score modificado.

\textbf{Caso de uso:} En un conjunto de salarios, si la media es muy superior a la mediana, indica presencia de unos pocos salarios muy altos (outliers). La mediana y MAD darán una mejor idea de la tendencia central y dispersión típica.

\subsection{Visualizaciones avanzadas}

\subsubsection{Distribuciones univariantes}
\begin{itemize}
    \item \textbf{Histograma con anchos de banda variables:} La elección del número de bins afecta la percepción. La regla de Freedman-Diaconis propone un ancho de bin óptimo:
    \[
    \text{ancho} = 2 \cdot \text{IQR} \cdot n^{-1/3}
    \]
    Útil para distribuciones con colas pesadas. \textit{Ejemplo:} Al histogramar ingresos, esta regla evita bins demasiado anchos que oculten la cola de altos ingresos.

    \item \textbf{KDE (Kernel Density Estimation):} Estima la función de densidad de probabilidad de forma continua:
    \[
    \hat{f}(x) = \frac{1}{nh}\sum_{i=1}^n K\left(\frac{x-x_i}{h}\right)
    \]
    donde $K$ es un kernel (usualmente gaussiano) y $h$ el ancho de banda. Un $h$ pequeño genera curvas rugosas; grande, sobre-suavizadas. \textit{Propósito:} Visualizar la forma de la distribución sin depender de bins discretos.

    \item \textbf{ECDF (Empirical Cumulative Distribution Function):}
    \[
    F_n(x) = \frac{1}{n}\sum_{i=1}^n \mathbf{1}_{x_i \leq x}
    \]
    Muestra la proporción de datos por debajo de un valor. Útil para comparar distribuciones y detectar desviaciones (ej. test de Kolmogorov-Smirnov).
\end{itemize}

\subsubsection{Relaciones bivariantes y multivariantes}
\begin{itemize}
    \item \textbf{Mapas de calor de correlación:} Matriz de correlación de Pearson (relación lineal):
    \[
    r = \frac{\sum (x_i - \bar{x})(y_i - \bar{y})}{\sqrt{\sum (x_i - \bar{x})^2 \sum (y_i - \bar{y})^2}}
    \]
    y Spearman (relación monótona, basada en rangos). \textit{Ejemplo:} En un dataset de viviendas, detectar alta correlación entre ``número de habitaciones'' y ``metros cuadrados'' para decidir si eliminar una.
\end{itemize}

\subsubsection{Datos temporales}
\begin{itemize}
    \item \textbf{Descomposición de series:} Modelo aditivo $Y_t = T_t + S_t + R_t$ o multiplicativo $Y_t = T_t \cdot S_t \cdot R_t$. Se estiman tendencia (media móvil), estacionalidad (promedios por período) y residuo.
    \item \textbf{Gráfico de autocorrelación (ACF):}
    \[
    ACF(k) = \frac{\sum_{t=k+1}^n (Y_t - \bar{Y})(Y_{t-k} - \bar{Y})}{\sum_{t=1}^n (Y_t - \bar{Y})^2}
    \]
    Mide la correlación de una serie consigo misma desplazada $k$ periodos. Útil para identificar dependencia temporal y órdenes en modelos ARIMA.
\end{itemize}

\subsection{Detección de outliers avanzada}

\subsubsection{Tipos de outliers}
\begin{itemize}
    \item \textbf{Globales:} Puntos que se desvían de todos los datos (ej. un ingreso de \$1M en un conjunto de ingresos medios).
    \item \textbf{Contextuales:} Anómalos en un contexto específico (ej. 30°C en invierno es normal en verano).
    \item \textbf{Colectivos:} Un grupo de puntos que juntos son anómalos (ej. un cluster de fraudes bancarios).
\end{itemize}

\subsubsection{Métodos univariados robustos}
\begin{itemize}
    \item \textbf{Rango intercuartílico (IQR):}
    \[
    \text{Límite inferior} = Q1 - 1.5 \times IQR,\quad \text{Límite superior} = Q3 + 1.5 \times IQR
    \]
    donde $IQR = Q3 - Q1$. \textit{Propósito:} Identificar outliers en cada variable de forma independiente.
    \item \textbf{Z-score modificado (basado en MAD):}
    \[
    M_i = \frac{0.6745(x_i - mediana(x))}{MAD}
    \]
    Umbral típico: $|M_i| > 3.5$. Funciona con distribuciones no normales y es resistente a outliers. \textit{Ejemplo:} Detectar transacciones financieras anómalas.
\end{itemize}

\subsubsection{Métodos multivariados}
\begin{itemize}
    \item \textbf{Distancia de Mahalanobis:}
    \[
    D_M(x) = \sqrt{(x - \mu)^T \Sigma^{-1} (x - \mu)}
    \]
    Mide distancia considerando correlaciones entre variables. Asume normalidad multivariante. Se considera outlier si $D_M^2 > \chi^2_{p,0.975}$. \textit{Aplicación:} Detectar perfiles de clientes atípicos en varias dimensiones (edad, ingresos, gastos).
    \item \textbf{Local Outlier Factor (LOF):} Compara la densidad local de un punto con la de sus $k$ vecinos. Se define:
    \[
    LOF(p) = \frac{\frac{1}{k}\sum_{o \in N_k(p)} lrd(o)}{lrd(p)}
    \]
    donde $lrd$ es la densidad local de alcance. Puntos con $LOF \gg 1$ son outliers. \textit{Ventaja:} Detecta outliers locales, no solo globales.
    \item \textbf{Isolation Forest:} Aísla puntos mediante particiones aleatorias. Los outliers requieren menos particiones para ser aislados. La puntuación de anomalía se basa en la profundidad media de aislamiento. \textit{Ejemplo:} Detectar fraudes en transacciones en tiempo real.
\end{itemize}

\subsection{Manejo de datos faltantes}

\subsubsection{Mecanismos de ausencia}
\begin{itemize}
    \item \textbf{MCAR (Missing Completely at Random):} La ausencia no depende de ninguna variable. Ej: error de transcripción aleatorio.
    \item \textbf{MAR (Missing at Random):} La ausencia depende de otras variables observadas. Ej: las mujeres responden menos sobre ingresos.
    \item \textbf{MNAR (Missing Not at Random):} La ausencia depende del valor faltante mismo. Ej: personas con altos ingresos no declaran.
\end{itemize}

\subsubsection{Técnicas de imputación}
\begin{itemize}
    \item \textbf{Imputación simple:}
    \begin{itemize}
        \item \textbf{Media/Mediana:} Reemplazar por $\bar{x}$ o mediana. Rápido pero subestima varianza y distorsiona correlaciones.
        \item \textbf{Regresión:} Predecir el valor faltante usando otras variables: $x_{faltante} = \beta_0 + \beta_1 y + \dots$.
        \item \textbf{kNN:} Tomar valor promedio de los $k$ vecinos más cercanos (distancia euclidiana). Preserva estructura local.
    \end{itemize}
    \item \textbf{Imputación múltiple (MICE):} Crear $m$ conjuntos imputados (generalmente 5-10) mediante modelos iterativos (regresión, árboles, etc.). Luego combinar resultados con las reglas de Rubin. La varianza total incorpora la incertidumbre de la imputación. \textit{Ejemplo:} En encuestas de salud, imputar ingresos faltantes usando edad, educación y región.
\end{itemize}

\section{Feature Engineering Avanzado}

\subsection{Transformaciones de variables}

\subsubsection{Normalización vs. estandarización}
\begin{itemize}
    \item \textbf{Normalización (Min-Max):} $x' = \frac{x - x_{min}}{x_{max} - x_{min}}$. Escala al rango $[0,1]$. Útil cuando se requiere rango fijo (ej. redes neuronales con activaciones sigmoide).
    \item \textbf{Estandarización (Z-score):} $x' = \frac{x - \mu}{\sigma}$. Media 0, desviación 1. Necesaria para modelos lineales, SVM, PCA.
\end{itemize}

\subsubsection{Transformaciones para asimetría}
\begin{itemize}
    \item \textbf{Logarítmica:} $x' = \log(x)$. Para datos positivos con cola larga derecha (ingresos, precios). Reduce asimetría y estabiliza varianza.
    \item \textbf{Box-Cox:}
    \[
    x' = 
    \begin{cases}
    \frac{x^\lambda - 1}{\lambda} & \lambda \neq 0 \\
    \log(x) & \lambda = 0
    \end{cases}
    \]
    Requiere $x > 0$. El parámetro $\lambda$ se estima maximizando la log-verosimilitud. \textit{Ejemplo:} Transformar área de viviendas para mejorar linealidad con el precio.
    \item \textbf{Yeo-Johnson:} Extensión que acepta valores negativos. Definición similar pero con modificaciones para $x$ negativos.
\end{itemize}

\subsubsection{Escalado robusto}
\[
x' = \frac{x - mediana(x)}{IQR}
\]
Basado en mediana e IQR, resistente a outliers. Útil cuando hay valores extremos que distorsionarían la estandarización clásica. \textit{Aplicación:} Preprocesar datos financieros con transacciones anómalas.

\subsection{Creación de nuevas características}

\subsubsection{Términos polinomiales e interacciones}
\begin{itemize}
    \item \textbf{Polinomiales:} $x^2, x^3$ capturan relaciones no lineales. Incluir hasta grado $d$ con precaución para evitar overfitting. \textit{Ejemplo:} Relación entre edad y riesgo de enfermedad puede ser cuadrática (U invertida).
    \item \textbf{Interacciones:} $x_i \cdot x_j$ capturan sinergias. Fundamental en modelos lineales para modelar efectos que no son aditivos. \textit{Ejemplo:} En predicción de precios, interacción entre ``número de habitaciones'' y ``metros cuadrados'' puede ser más informativa que las variables por separado.
\end{itemize}

\subsubsection{Binning y discretización}
\begin{itemize}
    \item \textbf{Igual ancho:} Dividir rango en $k$ intervalos de igual longitud. Sensible a outliers.
    \item \textbf{Igual frecuencia:} Intervalos con el mismo número de observaciones. Útil para distribuciones muy asimétricas.
    \item \textbf{Basado en árboles:} Usar los puntos de corte que un árbol de decisión encuentra para maximizar pureza. \textit{Ejemplo:} Convertir edad en grupos etarios para modelos de riesgo crediticio.
\end{itemize}

\subsubsection{Variables temporales}
\begin{itemize}
    \item \textbf{Rezagos (lags):} $y_{t-1}, y_{t-2}$ para series temporales. Permiten usar el pasado como predictor.
    \item \textbf{Ventanas móviles:} Media móvil, desviación móvil:
    \[
    MA_t = \frac{1}{k}\sum_{i=0}^{k-1} y_{t-i},\quad SD_t = \sqrt{\frac{1}{k-1}\sum_{i=0}^{k-1} (y_{t-i} - MA_t)^2}
    \]
    Capturan tendencias y volatilidad local.
    \item \textbf{Características cíclicas:} Para hora del día, día de semana, usar seno y coseno para preservar circularidad:
    \[
    \sin\left(\frac{2\pi \cdot \text{hora}}{24}\right),\quad \cos\left(\frac{2\pi \cdot \text{hora}}{24}\right)
    \]
    Evita el salto artificial entre 23h y 0h.
\end{itemize}

\subsubsection{Variables agregadas por grupo}
Por categoría (cliente, producto, región) se pueden calcular estadísticos como media, desviación, conteo, min, max. \textit{Ejemplo:} Gasto promedio por cliente en el último mes, frecuencia de compras, etc. Estas características resumen el comportamiento histórico.

\subsection{Codificación de variables categóricas}

\subsubsection{One-hot encoding}
Crear $k-1$ columnas binarias. Problema: alta dimensionalidad para muchas categorías. \textit{Usar cuando:} Número de categorías pequeño ($<10$).

\subsubsection{Target encoding}
Reemplazar la categoría por la media de la variable objetivo para esa categoría:
\[
x_{cat} = \frac{1}{n_{cat}} \sum_{i \in cat} y_i
\]
Requiere validación cruzada para evitar leakage. \textit{Útil para:} Categorías con muchas etiquetas y relación fuerte con el target. \textit{Ejemplo:} Códigos postales en predicción de precios de viviendas.

\subsubsection{Frequency encoding}
Reemplazar por frecuencia de aparición. Útil cuando la frecuencia es informativa (ej. palabras raras en NLP).

\subsubsection{Weight of Evidence (WoE)}
Para clasificación binaria:
\[
WoE = \ln\left(\frac{\% \text{eventos}}{\% \text{no eventos}}\right)
\]
Mide separación entre clases. Común en scoring crediticio.

\subsubsection{Embeddings para alta cardinalidad}
Aprender representaciones densas (ej. con redes neuronales) para variables como códigos postales o IDs de usuario. Se entrenan junto con el modelo principal.

\subsection{Selección de características}

\subsubsection{Métodos filtro}
\begin{itemize}
    \item \textbf{Correlación:} Eliminar variables con alta correlación ($>0.9$) para reducir multicolinealidad.
    \item \textbf{ANOVA / F-test:} Para clasificación, mide separabilidad entre clases. Se basa en la relación:
    \[
    F = \frac{\text{varianza entre grupos}}{\text{varianza dentro de grupos}}
    \]
    \item \textbf{Información mutua:}
    \[
    I(X;Y) = \sum_{x,y} p(x,y) \log\frac{p(x,y)}{p(x)p(y)}
    \]
    Captura cualquier tipo de dependencia (no solo lineal).
\end{itemize}

\subsubsection{Métodos wrapper}
\begin{itemize}
    \item \textbf{Selección hacia adelante (forward):} Añadir variable que más mejora métrica.
    \item \textbf{Eliminación hacia atrás (backward):} Eliminar variable que menos perjudica.
    Costosos computacionalmente, pero pueden encontrar interacciones.
\end{itemize}

\subsubsection{Métodos embeddeds}
\begin{itemize}
    \item \textbf{Importancia de árboles:} Basada en reducción de impureza (Gini o entropía) promediada en todos los árboles.
    \item \textbf{Regularización L1 (Lasso):}
    \[
    \min_{\beta} \sum (y - X\beta)^2 + \lambda \sum |\beta_j|
    \]
    Fuerza coeficientes a cero, seleccionando variables automáticamente. \textit{Ejemplo:} En regresión lineal con muchas características, Lasso identifica las más relevantes.
\end{itemize}

\section{Pipelines y automatización}

\subsection{Construcción de pipelines en scikit-learn}
Encadenar transformaciones y modelo final mediante `Pipeline`. Ventajas:
\begin{itemize}
    \item Evitar leakage (las transformaciones se ajustan solo en entrenamiento).
    \item Facilitar validación cruzada y búsqueda de hiperparámetros.
    \item Garantizar reproducibilidad.
\end{itemize}

\subsection{Feature engineering automático}
Herramientas como \texttt{Featuretools} implementan \textit{deep feature synthesis}: generan automáticamente características a partir de relaciones entre tablas (agregaciones, transformaciones). Útil para datos relacionales (clientes, transacciones, productos).

\section{Casos de uso integrados}

\subsection{Predicción de precios de viviendas}
\begin{itemize}
    \item \textbf{EDA:} Detectar asimetría en precio (aplicar log), outliers en área (usar IQR), correlaciones altas (eliminar redundantes). Visualizar relación entre ubicación y precio con boxplots.
    \item \textbf{Feature engineering:} Crear interacciones (habitaciones $\times$ baños), variables temporales (año de construcción transformado con Box-Cox), codificación de ubicación con target encoding, escalado robusto de características numéricas.
\end{itemize}

\subsection{Detección de fraude financiero}
\begin{itemize}
    \item \textbf{EDA:} Identificar outliers globales (transacciones con montos extremos) y contextuales (compras nocturnas inusuales). Analizar distribuciones por tipo de comercio.
    \item \textbf{Feature engineering:} Variables temporales (hora, día semana), agregadas por cliente (frecuencia, monto promedio, desviación en últimas 24h), distancia entre transacciones consecutivas, uso de Isolation Forest para generar puntuación de anomalía como nueva característica.
\end{itemize}

\section{Conclusión y conexión con modelado}
Un EDA profundo y un feature engineering cuidadoso no son pasos previos al modelado; son parte integral del proceso. Las decisiones tomadas aquí determinan la capacidad del modelo para generalizar. En las siguientes sesiones, se abordará cómo estos features alimentan modelos supervisados y no supervisados, y cómo la correcta aplicación de estas técnicas se refleja en métricas de rendimiento superiores.

\end{document}